\documentclass[11pt]{article}
\usepackage{amsmath,amssymb,amsthm}
\usepackage[margin=1in]{geometry}
\usepackage{hyperref}

\newtheorem{remark}{Remark}

\title{Period-doubling bifurcations of the logistic map:\\
\large deriving the exact polynomials with symbolic algebra}
\author{}
\date{\today}

\begin{document}
\maketitle

\begin{abstract}
The logistic map $x\mapsto rx(1-x)$ undergoes a cascade of period-doubling
bifurcations as the parameter $r$ increases from $1$ to $\approx 3.57$, after
which the dynamics become chaotic.  The parameter values at which each doubling
occurs are algebraic numbers --- they satisfy polynomial equations with integer
coefficients.  These notes show how to derive those polynomials from scratch
using nothing beyond the chain rule and polynomial long division, and explain
how the FORM symbolic-algebra program in this directory automates the
computation.  The ratios of successive bifurcation intervals converge to the
Feigenbaum constant $\delta\approx 4.669$, a universal quantity that appears in
every smooth one-hump map, not just the logistic one.
\end{abstract}

\tableofcontents

%======================================================================
\section{The logistic map}
\label{sec:logistic}
%======================================================================

Take a number $x_0\in[0,1]$ and iterate
\begin{equation}
\label{eq:logistic}
x_{n+1} = f_r(x_n) = r\,x_n(1-x_n),
\end{equation}
where $r>0$ is a parameter you get to choose.  This is the \emph{logistic map},
popularized by Robert May in 1976 as a toy model for population dynamics: $x_n$
is the fraction of some carrying capacity at generation $n$, and $r$ controls
growth.

What makes it interesting is how its long-term behavior changes as $r$
increases:
\begin{itemize}
\item $r<1$: every orbit converges to $x=0$ (the population dies).
\item $1<r<3$: every orbit converges to a single equilibrium value.
\item $3<r<1+\sqrt{6}\approx 3.449$: the equilibrium becomes unstable, and
      orbits settle into a period-2 oscillation --- bouncing forever between
      two values.
\item As $r$ increases further, the 2-cycle itself destabilizes and splits into
      a 4-cycle, then an 8-cycle, then 16, and so on.
\item At $r_\infty\approx 3.5699\ldots$\, the cascade accumulates and the map
      becomes chaotic.
\end{itemize}
Each of the transitions above happens at a \emph{specific} value of $r$, and
those values satisfy polynomial equations.  The goal of this project is to
derive those equations by symbolic computation.

%======================================================================
\section{Fixed points and stability}
\label{sec:fixed}
%======================================================================

A \emph{fixed point} of $f_r$ is a value $x^*$ that maps to itself:
$f_r(x^*)=x^*$.  For the logistic map this gives
\[
rx^*(1-x^*) = x^* \quad\Longrightarrow\quad x^*\bigl(r(1-x^*)-1\bigr)=0,
\]
so either $x^*=0$ or $x^*=1-1/r$.  Both are fixed points for any $r>1$, but
only one of them can be \emph{stable} --- attracting nearby orbits.

The stability criterion comes from linearizing: if we perturb $x^*$ by a small
amount $\varepsilon$, then after one step the perturbation is approximately
$f_r'(x^*)\,\varepsilon$.  The number $\lambda = f_r'(x^*)$ is called the
\emph{multiplier}.  If $|\lambda|<1$ the perturbation shrinks (stable); if
$|\lambda|>1$ it grows (unstable).  The boundary $|\lambda|=1$ is where the
qualitative behavior changes --- a \emph{bifurcation}.

For the nontrivial fixed point $x^*=1-1/r$:
\begin{equation}
\label{eq:mult1}
\lambda = f_r'(x^*) = r\,(1-2x^*) = r\,\Bigl(1-2+\frac{2}{r}\Bigr) = 2-r.
\end{equation}
This satisfies $|\lambda|<1$ precisely when $1<r<3$.  At $r=3$ the multiplier
hits $-1$, the fixed point becomes unstable, and a period-2 cycle is born.

The sign matters: $\lambda=-1$ means orbits \emph{overshoot} past the fixed
point on each step, alternating sides.  That alternation is exactly what turns
into a genuine 2-cycle once the fixed point can no longer pull orbits back.

%======================================================================
\section{Period doubling: from cycles to algebra}
\label{sec:doubling}
%======================================================================

A \emph{period-2 cycle} of $f_r$ is a pair $\{x_1,x_2\}$ with $f_r(x_1)=x_2$
and $f_r(x_2)=x_1$.  Equivalently, $x_1$ and $x_2$ are fixed points of the
second iterate $f_r^2=f_r\circ f_r$ that are \emph{not} fixed points of $f_r$
itself.

The second iterate $f_r^2(x)=f_r(f_r(x))$ is a degree-4 polynomial in $x$
(each composition squares the degree).  Its fixed-point equation $f_r^2(x)=x$
has degree 4, and two of the four roots are the already-known fixed points of
$f_r$.  Factoring those out leaves a quadratic whose roots are the period-2
points:
\begin{equation}
\label{eq:period2quad}
r^2 x^2 - r(r+1)\,x + (r+1) = 0, \qquad
x = \frac{(r+1)\pm\sqrt{(r+1)(r-3)}}{2r}.
\end{equation}
The discriminant $(r+1)(r-3)$ is positive only for $r>3$, confirming that the
2-cycle exists only past the first bifurcation.

Now the same stability question applies to the 2-cycle.  Its multiplier is the
derivative of $f_r^2$ evaluated at either cycle point, which the chain rule
factors as
\[
(f_r^2)'(x_1) = f_r'(x_1)\cdot f_r'(x_2).
\]
Using Vieta's formulas on the quadratic~\eqref{eq:period2quad} ($x_1+x_2 =
(r+1)/r$, $x_1 x_2 = (r+1)/r^2$) and expanding gives the multiplier
\begin{equation}
\label{eq:mult2}
(f_r^2)'(x_1) = -r^2+2r+4.
\end{equation}
Setting this equal to $-1$ (the threshold for the next period doubling):
\[
-r^2+2r+4 = -1 \quad\Longrightarrow\quad r^2-2r-5=0
\quad\Longrightarrow\quad r=1+\sqrt{6}\approx 3.449.
\]
So the exact parameter value at which the 2-cycle destabilizes and a 4-cycle is
born is $r=1+\sqrt{6}$, satisfying the degree-2 polynomial $r^2-2r-5=0$.  This
is the \emph{bifurcation polynomial} for the $2\to 4$ transition.

%======================================================================
\section{The general bifurcation conditions}
\label{sec:general}
%======================================================================

The pattern generalizes.  At the bifurcation from period $m$ to period $2m$,
two things happen simultaneously:
\begin{enumerate}
\item \textbf{The cycle exists:} there is a point $x$ of exact period $m$,
      meaning $f_r^m(x)=x$ but $f_r^k(x)\ne x$ for $0<k<m$.
\item \textbf{The cycle is marginally unstable:} its multiplier equals $-1$,
      i.e.\ $(f_r^m)'(x)=-1$.
\end{enumerate}
These are two polynomial equations in $x$ and $r$.  The first has degree $2^m$
in $x$ (each composition of $f_r$ doubles the degree); the second has degree
$2^m-1$.  As the examples above show, there is a lot of redundancy: the
fixed-point equation for $f_r^m$ includes all fixed points of $f_r^k$ for every
$k$ dividing $m$, and these must be factored out to isolate the genuinely
period-$m$ points.

Once the redundancy is removed, we have a system of two polynomial equations in
$(x,r)$, and the goal is to \textbf{eliminate $x$}, producing a single
polynomial in $r$ alone: the bifurcation polynomial.  Its roots are the values
of $r$ at which a period $m\to 2m$ doubling occurs.

In principle this is a textbook operation --- computing a \emph{resultant}.  In
practice the expressions grow enormously (degree $2^m$ in $x$), and doing it by
hand past the $2\to 4$ case is impractical.  This is where symbolic algebra
earns its keep.

%======================================================================
\section{Eliminating $x$: the reduction strategy}
\label{sec:elimination}
%======================================================================

The FORM program implements the elimination via \emph{iterated polynomial long
division} in $x$ over the coefficient ring $\mathbb{Z}[r]$.

The setup: after building the iterate $f_r^{2m}(x)$ symbolically and
differentiating, the program holds two polynomial conditions, which it labels
$P$ and $N$ (``positive'' and ``negative,'' for reasons explained shortly):
\begin{align}
P(x,r) &= (f_r^{2m})'(x) - 1 = 0, \label{eq:Pdef}\\
N(x,r) &= (f_r^m)'(x) + 1 = 0. \label{eq:Ndef}
\end{align}
Condition $N$ says the period-$m$ multiplier is $-1$ (the cycle loses
stability).  Condition $P$ says the period-$2m$ multiplier is $+1$ (the new
cycle is being born).  At the exact bifurcation value of $r$, both conditions
are satisfied by the same $x$ --- they describe the same event from two sides.

A third polynomial, the \emph{constraint} $C(x,r)=f_r^m(x)-x$ (with
lower-period factors divided out), pins down which $x$-values are on the
period-$m$ cycle.  The program first uses $C$ to reduce the degrees of $P$ and
$N$: from $C=0$, solve for the highest power of $x$ and substitute into $P$ and
$N$ wherever that power appears.  This brings their $x$-degrees down.

Then the \emph{ladder} phase kicks in: use $N$ to reduce $P$, then the reduced
$P$ to further reduce $N$, alternating like the Euclidean algorithm.  At each
step, known trivial factors ($r=\pm2$ and $r=4$, corresponding to degenerate
cases) are stripped.  The iteration continues until both polynomials have
minimal degree in $x$.

Finally, the \emph{resultant} step: extract the coefficient of each power of
$x$ from the two reduced polynomials, form cross-products to cancel the leading
$x$-term, and take the GCD.  What survives is a polynomial in $r$ alone --- the
bifurcation polynomial.

\begin{remark}[Why $r=\pm2$ and $r=4$ are trivial]
These values appear as roots at every bifurcation level and are not
``interesting'' new bifurcation values.  $r=2$ places the critical point $x=1/2$
at a fixed point (a superstable cycle); $r=4$ makes $f_r$ surjective on $[0,1]$
(conjugate to the full tent map, all cycles unstable); $r=-2$ lies outside the
physically meaningful range $[0,4]$ but satisfies the polynomial equations,
which do not know about interval constraints.
\end{remark}

%======================================================================
\section{The Feigenbaum constant}
\label{sec:feigenbaum}
%======================================================================

Let $r_n$ denote the parameter value at which the period-$2^{n-1}\to 2^n$
bifurcation occurs.  The first few are:
\begin{center}
\renewcommand{\arraystretch}{1.25}
\begin{tabular}{cll}
$n$ & $r_n$ & approximate value \\\hline
1 & $3$ & $3$ \\
2 & $1+\sqrt{6}$ & $3.44949\ldots$ \\
3 & (root of a degree-4 poly) & $3.54409\ldots$ \\
4 & (root of a degree-8 poly) & $3.56441\ldots$ \\
$\vdots$ & & \\
$\infty$ & (transcendental) & $3.56995\ldots$
\end{tabular}
\end{center}
The intervals $r_n-r_{n-1}$ shrink rapidly, and their successive ratios converge
to a constant:
\begin{equation}
\label{eq:feigenbaum}
\delta = \lim_{n\to\infty}\frac{r_n-r_{n-1}}{r_{n+1}-r_n} = 4.6692016\ldots
\end{equation}
This is the \textbf{Feigenbaum constant}.  Its remarkable property is
\emph{universality}: the same constant $\delta$ governs the period-doubling
cascade of \emph{every} smooth unimodal (one-hump) map, not just the logistic
one.  Replace $rx(1-x)$ with $r\sin(\pi x)$, or $r(1-(2x-1)^4)$, or any other
smooth function with a single maximum --- the bifurcation ratios still converge
to $4.669\ldots$

Mitchell Feigenbaum discovered this universality in 1975 using a pocket
calculator.  The explanation involves renormalization-group ideas: the
period-doubling operation acts on the space of unimodal maps, and $\delta$ is
the unstable eigenvalue of its linearization at a fixed point of that operator.
The algebraic approach taken here --- computing the exact bifurcation
polynomials --- provides complementary information: having the polynomials
(rather than floating-point roots) opens the door to studying $\delta$ through
the algebraic structure of the equations themselves.

%======================================================================
\section{What the code computes}
\label{sec:code}
%======================================================================

The FORM program \texttt{feigenbaum.frm} automates the full pipeline described
above.  A single call \texttt{\#call LogisticBifurcations(B,N)} computes the
bifurcation polynomial for the period $B\cdot 2^{N-1}\to B\cdot 2^N$
transition.  The steps are:

\begin{enumerate}
\item \textbf{Build the iterates.}  Starting from $x$, compose $f_r(x)=rx(1-x)$
      symbolically $B\cdot 2^N$ times.  At the halfway point (step $B\cdot 2^{N-1}$),
      snapshot $f_r^m(x)$ for the bifurcation conditions; at the quarter point,
      snapshot the lower-period polynomial to factor out later.

\item \textbf{Differentiate.}  Take $\partial/\partial x$ of both conditions to
      obtain $P$ and $N$ (equations~\eqref{eq:Pdef}--\eqref{eq:Ndef}).

\item \textbf{Remove lower-period solutions.}  Divide the period-$m$
      fixed-point equation by the period-$m/2$ equation (using FORM's built-in
      \texttt{div\_}), so only genuinely new cycle points remain.

\item \textbf{Reduce.}  Use the constraint to lower the $x$-degree of $P$ and
      $N$, then cross-reduce them against each other (the ``ladder''), stripping
      $r=\pm2,4$ at each step.

\item \textbf{Eliminate $x$.}  Extract $x$-coefficients from the reduced
      ladders, form cross-products, and take the GCD to obtain a polynomial in
      $r$ alone.

\item \textbf{Output.}  The result is stored in \texttt{RSolution($m$)} ---
      the minimal polynomial whose roots are the bifurcation parameter values.
\end{enumerate}

The main bottleneck is step~1: each composition doubles the degree, so
$f_r^{2m}$ has degree $2^{2m}$ in $x$.  The $1\to 2$ and $2\to 4$ bifurcations
are fast; higher doublings are expensive.

%======================================================================
\section{A hidden symmetry: conjugacy to $z^2+c$ and the variable $y=(r-1)^2$}
\label{sec:conjugacy}
%======================================================================

Something unexpected happens when you substitute $r = 1+t$ into the bifurcation
polynomials and expand: \emph{every odd power of $t$ vanishes}.  The $2\to4$
polynomial $r^2-2r-5$, for instance, becomes
\[
(1+t)^2 - 2(1+t) - 5 = t^2 - 6,
\]
so with $y=t^2=(r-1)^2$ the bifurcation condition is simply $y=6$.  The $3\to6$
polynomial (degree~6 in~$r$) likewise reduces to a cubic in~$y$, and the
$4\to 8$ polynomial (degree~12) to a sextic:
\begin{center}
\renewcommand{\arraystretch}{1.3}
\begin{tabular}{ccc}
bifurcation & $P(r)=0$ & $Q(y)=0$, \quad $y=(r-1)^2$ \\\hline
$2\to 4$ & $r^2-2r-5$ & $y-6$ \\[2pt]
$3\to 6$ & $r^6-6r^5+4r^4+24r^3$ & $y^3-11y^2+37y-108$\\
         & $\quad{}-14r^2-36r-81$ & \\[2pt]
$4\to 8$ & degree 12 & $y^6-18y^5+122y^4-516y^3$ \\
         &           & $\quad{}+1508y^2-1872y+5688$
\end{tabular}
\end{center}
The degree is halved every time: a degree-$2d$ polynomial in $r$ that is
symmetric about $r=1$ becomes a degree-$d$ polynomial in $y$.  But
\emph{why} should the bifurcation polynomials care about the point $r=1$?

\subsection{Finding the conjugacy}

Both the logistic map $f_r(x)=rx-rx^2$ and the ``standard'' quadratic
$g_c(z)=z^2+c$ are quadratics in their respective variables, so it is natural
to ask: is there an affine change of variables $z=\alpha x+\beta$ that turns
one into the other?

If such a change exists, then whenever $x_{n+1}=f_r(x_n)$ we need
$z_{n+1}=g_c(z_n)$, i.e.\
\[
\alpha\, f_r(x) + \beta = (\alpha x+\beta)^2 + c.
\]
The left side is $\alpha(rx-rx^2)+\beta = -\alpha r\,x^2 + \alpha r\,x+\beta$.
The right side is $\alpha^2 x^2+2\alpha\beta\,x+\beta^2+c$.  Matching
coefficients of $x^2$, $x$, and $1$:
\begin{align*}
x^2\!: & \quad -\alpha r = \alpha^2
        & \Longrightarrow\quad \alpha &= -r, \\
x^1\!: & \quad \alpha r = 2\alpha\beta
        & \Longrightarrow\quad \beta &= \tfrac{r}{2}, \\
x^0\!: & \quad \beta = \beta^2+c
        & \Longrightarrow\quad c &= \tfrac{r}{2}-\tfrac{r^2}{4}.
\end{align*}
So the unique affine conjugacy is
\begin{equation}
\label{eq:conjugacy}
z = -r x + \tfrac{r}{2} = r\!\left(\tfrac12 - x\right), \qquad
c = \frac{r}{2} - \frac{r^2}{4}.
\end{equation}
Every orbit of $f_r$ maps, point by point, to an orbit of $g_c$.  In
particular, cycles of $f_r$ correspond to cycles of $g_c$ of the same period,
with the same multiplier, so the two maps undergo bifurcations at exactly the
same parameter values.

\subsection{Why the symmetry follows}

The key is in the formula for $c$.  Complete the square:
\begin{equation}
\label{eq:c-of-y}
c = \frac{r}{2}-\frac{r^2}{4}
  = -\frac{r^2-2r}{4}
  = -\frac{(r-1)^2-1}{4}
  = \frac{1-(r-1)^2}{4}.
\end{equation}
Write $t=r-1$ and $y=t^2=(r-1)^2$:
\[
c = \frac{1-y}{4}.
\]
The parameter $c$ depends on $r$ \emph{only through $(r-1)^2$}.  Two values of
$r$ that sit symmetrically about $r=1$ --- say $r=1+t_0$ and $r=1-t_0$ ---
give the same $y=t_0^2$, hence the same $c$, hence \emph{exactly the same
quadratic map $g_c$}.  Every dynamical property (cycle existence, multiplier,
bifurcation threshold) is determined by $c$, so it must be invariant under
$r\mapsto 2-r$.

For the bifurcation polynomials this means: if $r_0$ is a root then so is
$2-r_0$.  But a polynomial whose roots come in $\{r_0,\,2-r_0\}$ pairs is
exactly one that is symmetric about $r=1$, i.e.\ $P(1+t)=P(1-t)$ --- which is
to say $P(1+t)$ contains only even powers of $t$.  Setting $y=t^2$ then
produces a polynomial in $y$ of half the degree, with no leftover square roots.

Setting $y=t^2$ then collapses each conjugate pair into a single root of a
polynomial $Q(y)$ of half the degree.  The relationship $y = 1-4c$ is linear,
so the $Q(y)$ are, up to an affine change, the bifurcation polynomials of the
\emph{Mandelbrot set} --- the polynomials whose roots are the $c$-values at
which $g_c$ undergoes period doubling.  For example, the $2\to4$ bifurcation
at $y=6$ corresponds to $c=(1-6)/4=-5/4$, which is the period-2 bifurcation
point of the Mandelbrot set; the $1\to2$ bifurcation at $r=3$ gives $y=4$ and
$c=-3/4$, the tip of the main cardioid.

\begin{remark}[Degree halving as a practical tool]
Working with $Q(y)$ instead of $P(r)$ is not just aesthetically cleaner --- it
halves the degree of every polynomial the computer must manipulate.  For the
$4\to 8$ bifurcation the saving is modest (degree~6 vs.~12), but at higher
levels the expressions grow exponentially and every factor of two in degree is a
significant win.  One could, in principle, run the entire elimination in the
$c$-parametrization (i.e., work with $g_c$ directly) and avoid the redundant
factor altogether.
\end{remark}

\begin{thebibliography}{9}
\bibitem{May1976}
R.~M.~May, \emph{Simple mathematical models with very complicated dynamics},
Nature \textbf{261} (1976), 459--467.
\bibitem{Feigenbaum1978}
M.~J.~Feigenbaum, \emph{Quantitative universality for a class of nonlinear
transformations}, J.~Stat.\ Phys.\ \textbf{19} (1978), 25--52.
\bibitem{Strogatz2015}
S.~H.~Strogatz, \emph{Nonlinear Dynamics and Chaos}, 2nd~ed., Westview Press, 2015.
\end{thebibliography}

\end{document}
